\section{S11: The Stray Longitudinal --- and Every Question It Raises Is One
Question}\label{step:S11}

\stagefield{Part / anchor}{Part I --- Light. Sector~1, step~3.}
\claimstatus{\StatusText{Computed Spectrum / Open Physical Status}}
\stagefield{Purpose}{Lift S10's longitudinal zero when the brane is allowed to
resist compression, enter the brane compression modulus as a postulated knob
with a named retirement condition, and state what the resulting spectrum does
and does not establish.}

\paragraph{Status.}
The longitudinal zero is lifted and its cone is computed. Its physical status
is not: whether the mode radiates, is observable, or constitutes a
Lorentz-breaking cone remains open because the required brane--bulk interface
coupling law has not been built.

\paragraph{Departure --- one statement, not three.}
\textbf{The second mode's entire physical status --- radiative, observable,
Lorentz-breaking, or none of the above --- reduces to one unbuilt object: the
brane--bulk interface coupling law that says whether and how matter couples to
it.} The law is linear and was deferred by choice, not by difficulty; it is the
work of S11b. The three questions carried by the walk are therefore one:
\begin{center}
\begin{tabular}{@{}L{0.60\textwidth}L{0.27\textwidth}@{}}
\toprule
question & reduces to \\
\midrule
does the longitudinal radiate into the bulk, or stay bound? & the coupling law \\
is light's confinement unconditional, or does it rest on a
polarisation-overlap argument? & the coupling law \\
does a second characteristic speed break Lorentz invariance for us? & the coupling law \\
\bottomrule
\end{tabular}
\end{center}

\paragraph{What S11 does not deliver.}
S11 does \textbf{not} decide bound versus radiating; does not establish that
light's confinement is unconditional; does not establish whether the
longitudinal is observable or unobservable; and supplies no leakage rate.

\paragraph{Why the extra mode is not a defect being repaired.}
Exact Maxwell would be the \textbf{failure}: Maxwell puts charge in by hand, so
a model matching it exactly would have no way to anchor charge physically. The
extra mode is that anchor; it is what made the drum-head charge picture click
(\StageFile{research/pde\_ledger\_v3/CHARTER.md\#falsification-standard}).
The token \texttt{FAIL\_CAUCHY\_STRAY\_LONGITUDINAL} is misnamed; its prefix is
not a verdict.

\paragraph{Identification, flagged before use.}
The field \(u\) is the material displacement of the stuff whose density is
\(\rho_{\mathrm{br}}\). S9's kinetic term
\(\tfrac12\rho_{\mathrm{br}}(\partial_tu)^2\) is kinetic energy only if
\(\dot u\) is the velocity of the material carrying that inertia. Linearised
continuity therefore applies,
\[
  \delta\rho_{\mathrm{br}}=-\rho_{\mathrm{br}}\nabla\!\cdot u,
\]
and the medium's equation of state makes compression cost energy. This is not
a modelling choice. It would have been wrong if \(u\) were an
orientation/director field: then \(\nabla\!\cdot u\) would be splay rather than
a density change, and the channel would need a separate argument.

\paragraph{Correction 1 --- the invariant count and the method.}
The derivative \(\partial_i u_j\) decomposes into trace, symmetric-traceless,
and antisymmetric parts. S10 did not forget a term: it retained one invariant
of several. Curl-only
stiffness charges twist and nothing else, so the pure-trace longitudinal wave
was free rather than forbidden. But the walk's claim of ``exactly three
coefficients for every \(D\geq2\)'' was \textbf{wrong} under proper rotations:
\[
  N_{SO}=\{D=2\mapsto4,\ D=3\mapsto3,\ D=4\mapsto4,\ D=5\mapsto3\},
  \qquad N_O=3\quad\hbox{for every }D.
\]
The method, not its arithmetic, was the error. Decomposing into pieces that do
not mix and counting them misses cross-pairings between isomorphic summands.
At \(D=2\), the trace and \(\Lambda^2\) are both \(SO(2)\) scalars; at \(D=4\),
\(\Lambda^2\) splits into self-dual and anti-self-dual pieces. The additional
reflection-odd invariants are
\[
  (\operatorname{tr}G)\,\epsilon^{ij}G_{ij}\quad(D=2),
  \qquad
  \epsilon_{ijkl}G_{ij}G_{kl}\quad(D=4).
\]
Five independent derivations agree on the corrected counts. The \(D=2\)
invariant is not a total derivative: its Euler--Lagrange operator is non-zero,
so an \(SO(2)\)-invariant Lagrangian can mix longitudinal and transverse
sectors. The \(D=4\) invariant is a total derivative and adds no operator
structure. Thus sector separation is a \textbf{\(D=3\) fact}, not a structural
one.

% Every rendered occurrence carries the status label required by the record.
% The underbrace is an annotation, not part of the symbol or an extra index.
\newcommand{\SXIcomp}{\underbrace{B_{\mathrm{comp}}}_{\text{\scriptsize postulated}}}

\paragraph{The spectrum.}
With the brane compression modulus \(\SXIcomp\), postulated rather than
derived, the trace term enters strictly as
\[
  \rho_{\mathrm{br}}\omega^2a
  =\mu_R\bigl[k^2a-k(k\!\cdot\!a)\bigr]
   +\SXIcomp\,k(k\!\cdot\!a).
\]
For \(k\!\cdot\!a=0\), the new term vanishes identically: compression cannot
touch light, and S9/S10 are not reopened. For \(a\parallel k\), the
\(\mu_R\) bracket vanishes and the old zero is replaced, not perturbed.

The two independently written engines compute
\[
  M_{ij}=\mu_R(k\!\cdot\!k)\delta_{ij}
         +(\SXIcomp-\mu_R)k_i k_j,
\]
and obtain:
\begin{center}
\begin{tabular}{@{}L{0.24\textwidth}L{0.65\textwidth}@{}}
\toprule
quantity & computed value \\
\midrule
transverse &
\(\omega^2=(\mu_R/\rho_{\mathrm{br}})k^2\), nullity \(D-1\), kernel
perpendicular to \(k\); unchanged from S9/S10 \\
longitudinal &
\(\omega^2=(\SXIcomp/\rho_{\mathrm{br}})k^2\), nullity \(1\), kernel
parallel to \(k\) \\
cross-sector &
the perpendicular root is independent of \(\SXIcomp\); the parallel root is
independent of \(\mu_R\); both computed residuals are zero \\
degeneracy & exactly \(\SXIcomp=\mu_R\) \\
\bottomrule
\end{tabular}
\end{center}

The dimensions close in general \(D\) as
\[
  [\rho_{\mathrm{br}}]=(-D,0,1),
  \qquad
  [\mu_R]=[\SXIcomp]=(2-D,-2,1),
  \qquad
  [c_\gamma]=[c_L]=(1,-1,0).
\]
Thus the longitudinal speed is the derived quantity
\[
  c_L^2=\frac{\SXIcomp}{\rho_{\mathrm{br}}}.
\]
Every computed value agrees with the preregistration
(\StageFile{research/pde\_ledger\_v3/steps/S11\_PREREGISTERED\_PREDICTION.md},
commit \texttt{67d919bd}) except the move-2 invariant count: there the
preregistration was wrong and the engines were right.

\paragraph{Controls.}
The form control replaces the trace invariant with the
symmetric-traceless one and moves both roots:
\[
  \omega_\perp^2=\frac{\mu_R+\mu_{\mathrm{br}}}{\rho_{\mathrm{br}}}k^2,
  \qquad
  \omega_\parallel^2=\frac{4\mu_{\mathrm{br}}}{3\rho_{\mathrm{br}}}k^2.
\]
The transverse root thereby acquires \(\mu_{\mathrm{br}}\), showing that the
trace invariant is the unique one that lifts the longitudinal while leaving
light untouched. The coefficient control merely rescales \(\SXIcomp\) and
moves only the parallel root, so it cannot test that shape claim. The form
control's \(4/3\) independently reproduces the corpus's rejected-Cauchy-branch
coefficient from an engine blind to that document.

\paragraph{Correction to the compression mechanism.}
The walk's first mechanism was \textbf{wrong}: it had brane material
``squeezing out into the bulk,'' which smuggled in a phase transition because
the bulk is the disordered phase. Compression does not make the brane
unordered. The correct mechanism is that the brane \textbf{thickens}: in-plane
compression raises areal density, accommodated either as higher four-dimensional
density, charged by the equation of state, or as a wider wall, charged by the
wall potential. The order parameter never changes; the sheet bulges into
\(\pm w\). The two channels take additive shares of one strain, so they are
springs in series, their compliances add, and the postulated \(\SXIcomp\) is
softer than either channel alone.

\paragraph{The flat direction and the named retirement condition.}
S7 records that the double well selects no slab width. If that flatness
survived, \(B_{\mathrm{wall}}=0\), hence the postulated \(\SXIcomp=0\), and
S10's zero would not lift. Width gradients lift it: a wave modulates the width,
tilts and stretches the interfaces, and costs
\(\sigma_{\mathrm{wall}}|\nabla W|^2\). The mode is therefore flat at \(k=0\)
and stiff as \(k^2\), predicting a flexural crossover with
\(\omega\propto k^2\) at long wavelength.

The cone above was computed with wall width \textbf{frozen}; unfreezing it
should soften the mode. If it does not, move~5 was wrong and must not be
reconciled away. The wall tension \(\sigma_{\mathrm{wall}}\) is S6-derived, so
the named retirement condition is
\texttt{DEFECT\_REGISTER\#c12}: the postulated \(\SXIcomp\) is to be retired
visibly at S6 if that derivation supplies it rather than leaving an independent
knob. The present knob count is an upper bound that can only improve.

\paragraph{Correction 2 --- bulk matching was overstated in both directions.}
Phase matching gives
\[
  k_w^2=k^2\left(\frac{c_L^2}{c_{s0}^2}-1\right).
\]
The walk read \(c_L>c_{s0}\) as ``radiates'' and \(c_L<c_{s0}\) as ``bound.''
\textbf{Neither direction is established.} Phase matching is kinematic only:
it determines whether a propagating bulk channel exists, not whether the mode
uses it, and not whether an evanescent solution forms a bound eigenmode. Both
claims require the absent interface coupling law. The omitted third case is
\(k_w=0\), grazing.

\paragraph{What is new.}
\begin{center}
\begin{tabular}{@{}lll@{}}
\toprule
item & class & why \\
\midrule
\(Q.\mathrm{brane}.\SXIcomp\) & \textbf{postulated} & retire under
\texttt{DEFECT\_REGISTER\#c12}; knob count is an upper bound \\
\(Q.\mathrm{brane}.c_L\) & \textbf{derived} & \(R5\), from the postulated
\(\SXIcomp\) and \(\rho_{\mathrm{br}}\) \\
the mode itself & \textbf{derived} & nullity of the dynamical matrix \\
\bottomrule
\end{tabular}
\end{center}

\noindent
The postulated \(\SXIcomp\) is a \textbf{brane} compression modulus. It is not
\(K_{\mathrm{br}}\), the rejected elastic branch's bulk modulus; not
\(B_{\mathrm{eff}}=\rho_{B0}^2/\chi_c\); and not the bulk medium's modulus. It
is registered with \texttt{aliases: []} and no borrowed loci. The registry
changes are ambient \(10\mathbin{\to}12\) and residue \(6\mathbin{\to}7\), with
residue
\(\{\hbar,m,K,\rho_0,\rho_{\mathrm{br}},\mu_R,\SXIcomp\}\). The discrete
payload \(\{n_{\mathrm{eos}}=5,D_{\mathrm{brane}}=3\}\) is unchanged and
remains off the continuous axis.

\paragraph{Inputs, by locus.}
S9 supplies the material-displacement kinetic term and the postulated
\(\rho_{\mathrm{br}}\) and \(\mu_R\); S10 supplies the longitudinal zero and
the in-plane identification. The equation of state supplies the compression
cost. S7 supplies the flat slab-width statement, and the S6-derived wall
tension supplies the proposed retirement route. The interface law is not an
input because it has not been built.

\paragraph{Regime.}
Linear response about the brane, with wall width frozen for the computed cone.
The predicted long-wavelength softening belongs to the width-unfrozen problem
and is not computed here.

\paragraph{What this settles about the light sector as a whole.}
First, \(c_L\) is \textbf{not} a gravitational wave. Gravity in this model is
the flow between draining defects, carried by flow plus Bernoulli pressure and
not by ripples or radiation
(\StageFile{docs/conceptual\_foundation.md}, line~348). The longitudinal is an
in-plane displacement of brane material. Separately, the corpus has no
mechanistic account of what a gravitational wave is.

Second, a second cone is a departure only if matter couples to it. The
transverse wave equation is Lorentz-invariant with invariant speed
\(c_\gamma\), and matter in this model is built from those modes: a throat is
held open by a trapped brane-shear standing wave. This is not circular; the
wave equation supplies the invariance, so a standing wave built from it must
restructure under boost, with its cost diverging as \(v\to c_\gamma\). The
walk's statement that a second cone by itself means observers see no single
invariant speed was wrong.

Third, uniform flow is unobservable while gradients are observable. If uniform
flow were observable, absolute motion would be detected; if gradients were
not, there would be no gravity. ``The bulk flow is invisible because we are
made of it'' and ``density gradients bend light'' are statements about
different derivatives.

Fourth, light gravitates but does not dissipate, and the second fact follows
from S9. Light carries stress-energy as a co-moving disturbance without loss,
while photons from distant galaxies show that light must not leak. A photon
can lose energy only into an available mode, and the bulk has no transverse
mode. Photon stability over cosmological distance is therefore a consequence
of bulk shear-freeness. The genuine loss channel is mode conversion where the
medium is inhomogeneous, near matter and not in vacuum.

% NOT \stagefield{Verification}: that macro suppresses Verification in the
% default build.  The limitations and the meaning of PASS must stay visible.
\paragraph{Verification.}
Two independently written engines agree on every computed value. The blind
Mathematica audit
(\StageFile{research/pde\_ledger\_v3/mathematica/S11\_stray\_longitudinal\_mathematica\_audit.wl})
was written first, imports nothing, and reads no registry. The SymPy audit
(\StageFile{research/pde\_ledger\_v3/scripts/S11\_stray\_longitudinal\_sympy\_audit.py})
was written while the Mathematica file was quarantined, then compared with its
byte-identical committed blob \texttt{55620a7f}; the builder's git use was
checked as \texttt{status} and \texttt{diff} only.

The engines initially disagreed on one valuable point. Mathematica asserted a
conclusion about transverse coupling while SymPy refused to conclude it; an
independent review had already found the Mathematica line unconditional and
unearned. Mathematica was repaired to stop making the claim, not steered to
match SymPy. It then diagnosed its own circularity and named the missing input.
There were eight review legs: two for build directives, two per engine, and two
per repair.

A demonstrated control hole was also closed. Replacing \(R5\) by
\(c_L-\sqrt{\mu_R/\rho_{\mathrm{br}}}\), the exact wrong relation this step
exists to distinguish, had left all five gates green. A new assertion now
substitutes each derived root into the registry residual that records it; four
distinct corruptions were ablation-checked, and \(R4\) is covered as well
(\texttt{DEFECT\_REGISTER\#f-r5}).

The reported gates are acceptance \StatusText{match}
\((12\mathbin{\to}7,\ 12\mathbin{\to}7,\ 12\mathbin{\to}6)\);
dimensional homogeneity \StatusText{pass}, including homogeneous \(R5\) with
dimension \((1,-1,0)\) and zero undetermined; able-to-fail \StatusText{pass};
and 11 tests. Both engines print \texttt{VERDICT: PASS}. \textbf{Those strings
are verdicts only on the engines' internal checks; they are not verdicts on the
physics, the absent interface coupling, or the mode's physical status.}

\paragraph{Known limits --- recorded, not fixed.}
\begin{enumerate}[leftmargin=2.2em]
\item Two of the six dimensional checks in the Mathematica file are
definitional identities, \((X-2Y)+2Y=X\), and cannot fail. Only the
\(\rho_{\mathrm{br}}\) check obtained an independent route. The guard as a
whole is able to fail and four independent ablations were caught, but those
two entries are \textbf{not counted as coverage}.

\item \texttt{ASSERTION\_12} in the SymPy file defines
\(c:=\sqrt{X}\) and then asserts \(c^2-X=0\). It is identically zero by
construction, cannot fail, and is \textbf{not counted as coverage}; the new
registry control supersedes it in substance.

\item The Mathematica file disclaims scope on the transverse tag but not the
parallel tag, although both arise from the same substitution with the same
absent operator. The scope limit applies to both channels.

\item The SymPy script was repaired while untracked, so there is no committed
pre-repair baseline. A reviewer's \texttt{/tmp} copy and a full hand
re-derivation of every load-bearing value substituted for one. This violated
the rule to commit before anything destructive.

\item A review leg's search for \texttt{K\_br} incidentally returned two lines
of \StageFile{research/pde\_ledger\_v3/V3\_STEP\_PLAN.md}, which was on its
do-not-read list. The file was not opened, and all reported results had already
been derived.
\end{enumerate}
